\section{Introduction}\label{sec:intro}

Weil's criterion \cite{Weil1952} states that the Riemann hypothesis is equivalent to the positivity of an explicit
Hermitian form $\Q$ on compactly supported test functions on the multiplicative group. The form has three parts, an
archimedean one, a pole term and a sum over primes, and positivity is known only for tests supported in a short window:
Yoshida \cite{YOSHIDA-HERMITIAN-1992} and Bombieri \cite{BOMBIERI-WEILQF-2000} for small supports, Connes and Consani
\cite{CC-WEILPOS-2020} for the archimedean place, and the finite window matrices of Connes, Consani and Moscovici (CCM)
\cite{CCM-ZST-2025}, whose bottoms are upper bounds for the infimum of $\Q$ on a window and whose positivity for all windows
is again RH. Bombieri already observed that the number of negative eigenvalues of finite truncations counts the zeros off the
line \cite{BOMBIERI-WEILQF-2000}; the spectral approach thus relocates the problem without reducing it.

This paper is about the shape of that problem. We fix one test function: Riemann's function $\Phi$, whose Fourier transform is
$\Xi$ itself \cite{Riemann1859}, in the normalisation of \cite[Lemma 7.1]{CCM-ZST-2025}, where its scaled versions serve as
the ``educated guess'' for the bottom vectors of the window matrices. Normalised, $f_0=\Phi/\|\Phi\|_2$ is strictly positive and
lies in the radical of $\Q$ together with all its translates and all its convolutions (\Cref{prop:radical}); this is the reason
the compactly supported class, on which the form is positive definite under RH \cite{Suzuki2025}, has to be enlarged.
Dividing every test by $f_0$ turns the Weil form into a Dirichlet form with a signed jump kernel (\Cref{thm:GS}):
\[
 \Q(f_0s)=\frac12\iint f_0(x)f_0(x')\,|s(x)-s(x')|^2\,d\nu(|x-x'|)\,dx\,dx',
\]
where the measure $\nu$ is completely explicit: a continuous density that is positive for $|x-x'|<\log\pl$, with $\pl$ the plastic
number, negative beyond, and positive atoms $\Lam(n)/\sqrt n$ at the prime-power distances $\log n$. The signed distributional kernel itself is
not new: Suzuki writes the Weil form as $\iint(-g'')(x-y)\,v(y)v(x)$ with the same prime atoms and an archimedean density with a
$\tfrac12\,\mathrm{Pf}(1/|t|)$ singularity \cite[\S2.5]{Suzuki2026}, \cite[\S3.5]{Suzuki2025}, and the statement ``$g$ is a screw
function in the sense of Krein--Langer if and only if RH holds'' \cite{Suzuki2026} is the conditional-positivity form of Weil's
criterion. What is new here is the substitution $g=f_0s$: it is the ground-state representation of Frank and Seiringer
\cite{FrankSeiringer2008}, applied to a signed L\'evy-type measure with a radical zero mode instead of a genuine ground state, and
it produces the closed form of the density, its sign change at the plastic number, and the differences $|s(x)-s(x')|^2$ that the
obstruction theorem needs. The Riemann hypothesis is then the statement that this form is nonnegative for every compact smooth $s$
(\Cref{cor:DOM}): one inequality between the primes and a negative density, with no spectral unknowns.

The main theorem is negative. The only unconditional positivity mechanism known for the Weil form is nonlocal: compression to
Sonin's space gives, for short supports, a minorant $W_\infty(g*g^*)\ge\operatorname{Tr}(\vartheta(g)S\vartheta(g)^*)$
\cite[Thm.~7.1]{Connes2026}, \cite{CC-WEILPOS-2020}. A natural way to go further is a local certificate: an identity or lower bound by a
manifestly nonnegative expression built from finite difference stencils with nonnegative weights, as one certifies positive
semidefinite matrices with zero row sums on finite graphs. \Cref{thm:nominorant} shows that every such minorant of $\Q(f_0\,\cdot)$
is zero, with no restriction on the size, diameter or number of stencils, and \Cref{cor:noCSS} that no sum-of-squares representation
exists. The mechanism is the radical: compact truncations of the translated canonical test have vanishing form value, and finitely many
translates of a positive integrable function are linearly independent, so every stencil would have to annihilate all sample values.
The obstruction is not a sign statement (\Cref{ex:control}); it holds for any positive integrable function in the radical of any functional
with the cutoff property, and the Weil form is a corollary. A simple companion statement excludes uniform margins on dense families (\Cref{prop:trade}).

The remaining sections record what the finite windows do and do not show (\Cref{sec:windows}), including an exact second-jet formula
for the prolate trial of \cite{CCM-ZST-2025} and the agreement of certified window bottoms with the decay law of \cite{XUEFENGZHU-2026}.
They also record a ledger of excluded proof shapes with the theorem that excludes each (\Cref{sec:ledger}), the precise open
problem and the two representations it leaves open, and the provenance of every statement: which statements are machine-checked
in Lean~4 (the core of the obstruction theorem is) and how automated reasoning tools were used (\Cref{sec:disclosure}). Nothing in this paper is a step toward RH in the sense of a new
inequality; its content is the exact location of the difficulty and the proof that a class of attacks on it cannot succeed.

\paragraph{Priority.} The test function and $\widehat\Phi=\Xi$ are Riemann's \cite{Riemann1859}, in the form of
\cite[(7.1)--(7.2), Lemma 7.1]{CCM-ZST-2025}; the same logarithmic-variable expansion appears in the preprint \cite{Freedman2026}.
The Weil form as a quadratic functional with finite truncations, their non-degeneracy for simple zeros (his Lemma~10), and the
question of linear independence of the functions $x^{-\rho}$, $\rho$ ranging over the zeros, are Bombieri's \cite{BOMBIERI-WEILQF-2000}; the window matrices and the
prolate trial are from \cite{CCM-ZST-2025}, with finite dictionaries in the preprints \cite{Groskin2026,Freedman2026}; the signed
distributional kernel and the screw-function form of the criterion are Suzuki's \cite{Suzuki2025,Suzuki2026}; positive operators
attached to the Weil distribution appear in \cite{Li2024}. To our knowledge the ground-state substitution \eqref{eq:GS} with its
explicit density and sign change, the description of the radical on the enlarged class, the obstruction \Cref{thm:nominorant} with its
corollaries, and \Cref{prop:trade} have not appeared before; we would be grateful for references.
